\chapter*{图片致谢}

本书的插图为原创。照片依各自的自由许可使用，谨向其作者致谢：

\begin{itemize}
  \item 第~\ref{ch:b2:microbial-metabolism}~章：谢尔盖·维诺格拉茨基的照片，1890 年代——公有领域（维基共享资源）。
  \item 第~\ref{ch:b2:genome-diversification}~章：芭芭拉·麦克林托克在其实验室中的照片，史密森学会 / 科学服务社，1947 年——公有领域（维基共享资源）。
  \item 第~\ref{ch:b2:blood-circulation}~章：威廉·哈维的肖像，以及其著作 \emph{Exercitatio anatomica de motu cordis}（1628）中前臂静脉的插图，Wellcome
        Collection——CC~BY~4.0（维基共享资源）；一个红细胞、一个血小板和一个淋巴细胞的扫描电子显微照片，Electron
        Microscopy Facility, National Cancer Institute——公有领域（维基共享资源）。
  \item 第~\ref{ch:b2:speciation}~章：John Gould 绘制的达尔文地雀图版，取自 \emph{Journal of Researches}（1845）——公有领域（维基共享资源）。
  \item 第~\ref{ch:b2:biogeochemical-cycles}~章：莫纳罗亚天文台（Forrest M. Mims III，2006 年）以及在该处采集气瓶样本的情景（Commander John Bortniak，1982 年），NOAA Photo Library
——公有领域（维基共享资源）。
\end{itemize}

每张照片的完整来源链接和许可证网址，均列于本书代码仓库中的 \texttt{images/book4/CREDITS.md}。

\bigskip
除照片与原创插图外，本卷各章中约七十幅图为 AI 生成插图，由 OpenAI 图像生成功能依据本书编辑撰写的提示词制作，并在纳入前经过生物学准确性审阅。每幅生成图像的完整提示词均记录于仓库中的 \texttt{images/book4/ai/PROMPTS.md}。
\clearpage
